%==============================================================================
% Chapitre 4 - Gravitation
%==============================================================================

\section{Introduction}

La gravitation est l'un des phénomènes physiques les plus familiers.

C'est elle qui maintient les objets à la surface de la Terre et qui est
responsable de la chute des corps.

En balistique, la gravité constitue l'une des principales forces agissant sur
un projectile après sa sortie du canon.

Sans gravité, un projectile poursuivrait son mouvement selon sa direction
initiale.

\begin{aretenir}

La gravité est responsable de la chute des projectiles.

Elle agit en permanence pendant toute la durée du vol.

\end{aretenir}

\section{Observation intuitive}

Lorsqu'un objet est lâché sans vitesse initiale, il tombe vers le sol.

Cette observation est connue depuis l'Antiquité.

Une expérience de pensée permet de comprendre l'influence de la gravité sur
les trajectoires.

Imaginons deux projectiles :

\begin{itemize}
\item le premier est simplement lâché ;
\item le second est tiré horizontalement.
\end{itemize}

Tous deux subissent la même accélération gravitationnelle.

Ils touchent donc le sol au même instant si l'on néglige la résistance de
l'air.

\section{La chute libre}

On parle de chute libre lorsqu'un objet est soumis uniquement à la gravité.

Dans ce cas idéal :

\begin{itemize}
\item aucun frottement n'est pris en compte ;
\item aucune autre force n'agit sur l'objet.
\end{itemize}

La chute libre constitue un modèle simplifié extrêmement utile pour comprendre
les trajectoires.

\section{Accélération gravitationnelle}

À proximité de la surface terrestre, tous les objets subissent
approximativement la même accélération gravitationnelle.

Cette accélération est notée $g$.

Sa valeur moyenne est :

\begin{equation}
g \approx 9,81~\mathrm{m/s^2}
\label{eq:gravitation_g}
\end{equation}

Cela signifie que la vitesse verticale d'un objet augmente d'environ
9,81~m/s chaque seconde.

\begin{exemple}

Après une seconde de chute libre :

\begin{itemize}
\item vitesse verticale : environ 9,81~m/s ;
\end{itemize}

Après deux secondes :

\begin{itemize}
\item vitesse verticale : environ 19,62~m/s.
\end{itemize}

\end{exemple}

\section{Poids et masse}

La masse et le poids sont souvent confondus.

La masse représente la quantité de matière contenue dans un objet.

Elle s'exprime en kilogrammes.

Le poids correspond à la force exercée par la gravité sur cet objet.

\begin{rigueur}

Le poids est donné par :

\begin{equation}
P = m,g
\label{eq:poids}
\end{equation}

où :

\begin{itemize}
\item $P$ représente le poids ;
\item $m$ représente la masse ;
\item $g$ représente l'accélération gravitationnelle.
\end{itemize}

\end{rigueur}

\section{Influence sur les trajectoires}

La gravité agit continuellement sur les projectiles.

Même lorsqu'une arme est tirée parfaitement à l'horizontale, le projectile
commence immédiatement à perdre de l'altitude.

Cette chute existe dès la sortie du canon.

À courte distance, elle est faible.

À longue distance, elle devient l'un des phénomènes dominants.

\begin{exemple}

Une balle de 7,62~mm quittant le canon à 850~m/s ne parcourt pas une ligne
droite.

Sa trajectoire s'incurve progressivement vers le sol sous l'effet de la
gravité.

\end{exemple}

\section{Trajectoire idéale}

Dans un modèle simplifié :

\begin{itemize}
\item la gravité est prise en compte ;
\item la résistance de l'air est négligée.
\end{itemize}

La trajectoire obtenue est une parabole.

Ce modèle est très utile pour comprendre les principes fondamentaux.

Cependant, les trajectoires réelles s'écartent progressivement de cette forme
à cause de la traînée aérodynamique.

\section{Influence de l'altitude}

La valeur de $g$ n'est pas parfaitement constante.

Elle varie légèrement :

\begin{itemize}
\item avec l'altitude ;
\item avec la latitude ;
\item avec la forme de la Terre.
\end{itemize}

Pour les armes légères, ces variations sont généralement négligeables.

Elles deviennent toutefois importantes dans certains domaines :

\begin{itemize}
\item artillerie ;
\item missiles ;
\item applications spatiales.
\end{itemize}

\section{Application à la balistique}

La gravité intervient dans :

\begin{itemize}
\item la chute du projectile ;
\item le réglage des organes de visée ;
\item la détermination des points sécants ;
\item le calcul des tables balistiques ;
\item les modèles numériques.
\end{itemize}

Elle constitue généralement le premier phénomène physique intégré dans un
modèle balistique.

\begin{simulation}

Le modèle balistique le plus simple consiste à appliquer uniquement la
gravité.

Même si ce modèle est très éloigné de la réalité, il permet déjà d'obtenir
une trajectoire approchée.

La plupart des moteurs balistiques commencent par ce modèle avant d'ajouter
progressivement :

\begin{itemize}
\item la traînée ;
\item le vent ;
\item la rotation terrestre ;
\item les effets gyroscopiques.
\end{itemize}

\end{simulation}

\section{À retenir}

\begin{aretenir}

\begin{itemize}
\item La gravité agit en permanence sur le projectile.
\item L'accélération gravitationnelle moyenne vaut environ 9,81~m/s$^2$.
\item La masse et le poids sont deux notions différentes.
\item La gravité provoque la chute des projectiles.
\item Sans résistance de l'air, la trajectoire est parabolique.
\end{itemize}

\end{aretenir}

